\documentclass[12pt]{article}
\usepackage{graphicx} % Required for inserting images
\usepackage{amsfonts}
\usepackage[a4paper, margin=1in]{geometry}
\usepackage[x11names, svgnames, dvipsnames]{xcolor}
\usepackage[normalem]{ulem}
\definecolor{Burgundy}{RGB}{144,0,32}


\title{ \color{Burgundy} Cover letter}
\date{}

\begin{document}
\maketitle
\pagenumbering{gobble}


\noindent Dear Editors,\\

We are submitting for your consideration our manuscript entitled “Quantum spin ice in three-dimensional Rydberg atom arrays”.\\


A quantum spin liquid (QSL) is an exotic state of matter having emergent excitations with fractional quantum numbers and is described by the deconfined phases of emergent gauge theories. Despite being proposed over fifty years ago by Anderson, the detection of QSLs in traditional condensed matter systems has proven elusive due to the absence of a local order parameter that can distinguish them from ordered states. \\



Recently there has been great interest in the quantum simulation of exotic states of matter using the Rydberg atom array platform. This platform provides a high degree of control over the parameters of the system along with the ability to measure non-local operators. For example, a recent experiment performed on this platform, [G. Semeghini et al., Science \textbf{374}, 1242 (2021)], observed promising signs of a $\mathbb{Z}_2$ spin liquid, so named because it is described by the deconfined phase of a $\mathbb{Z}_2$ gauge theory at low energies.\\


Is it possible then to realize on this platform a system where the more familiar $U(1)$ gauge theory (the theory of quantum electrodynamics) emerges at low energies? If it is, then by Polyakov’s argument, such a system would necessarily be 3+1 dimensional. In our manuscript, we argue that the answer is yes. We start with the Hamiltonian of a three-dimensional Rydberg atom array that includes long-range interactions between Rydberg atoms and study its phase diagram. We find that in an experimentally accessible regime, a $U(1)$ QSL can be realized. We then go on to describe detailed protocols for the measurement of various (often non-local) correlators which provide smoking gun signatures of the QSL state. We find that by tuning experimentally relevant parameters, one can access the confined, deconfined, and Higgs phases of the $U(1)$ gauge theory in three spatial dimensions. Consequently one can also access the deconfinement-confinement and the Higgs phase transitions.\\


We believe that our manuscript satisfies the following criteria for publication in PRX listed on the APS website:\\

\noindent \textbf{(1) “Establish a fruitful analogy or connection between different subfields or topical areas of physics, or between physics and other scientific disciplines.”}\\ 
\indent Our work is a proposal for the quantum simulation on a Rydberg atom array of a QSL described by a $U(1)$ gauge theory in 3+1 dimensions. Further, we show that both deconfinement-confinement and the Higgs transitions can be accessed in our model which are of interest in both high-energy physics and condensed matter physics. Therefore, our work is a fruitful connection between condensed matter physics, AMO (atomic, molecular, optical) physics, quantum information science, and high-energy physics.\\


\noindent \textbf{(2) “Push an established field into a new direction.”}\\
\indent	 Our work pushes the search of a $U(1)$ QSL, which has so far remained restricted to traditional condensed matter systems, in the direction of Rydberg atom arrays. With this comes a completely new set of experimentally relevant questions for QSLs related to non-equilibrium dynamics, lattice engineering, and detection schemes that are specific to this platform. We believe that these questions will lead to fruitful scientific developments in the future.\\

\noindent \textbf{(3) “Significantly advance the state of the art of a field.”}\\
\indent	Previous works have studied the quantum simulation of a gapped $\mathbb{Z}_2$ QSL in a two-dimensional Rydberg atom array. We advance the state of the art by focusing on a gapless $U(1)$ QSL which cannot be realized in two dimensions.\\

For the above reasons, we believe that our manuscript would particularly suit the interdisciplinary audience of PRX. We thank you for your consideration.\\

\noindent Yours sincerely,\\
Jeet Shah, Gautam Nambiar, Alexey V. Gorshkov, and Victor Galitski





\end{document}
